\chapter*{Conclusion} % chapter* je necislovana kapitola
\addcontentsline{toc}{chapter}{Conclusion} % rucne pridanie do obsahu
\markboth{CONCLUSION}{CONCLUSION} % vyriesenie hlaviciek
% See file \verb'zaver.tex' for information about recommended contents of the Conclusion chapter.

In Chapter \ref{ch:polehunter}, we developed a generator of cubic multipoles \texttt{polehunter} with minimum girths $3$, $4$ and $5$. This algorithm is significantly faster than other known generators \texttt{multigraph} and \texttt{pregraphs}, and by a theoretical background, we show a possibility to extend the algorithm to minimal girths $6$ and $7$.

In Chapter \ref{ch:flow_pairs}, we proved that graphs of oddness $2$ have a $1/2$-flow-pair and also graphs with a $1/2$-flow-pair admit a rich NZ $7$-flow. Hence, we used the $1/2$-flow-pair as the ``structure-in-the-middle''. We believe that the $1/2$-flow-pair can be used to prove more results in the NZ flow theory. For example, when someone proves the existence of the $1/2$-flow-pair for some uncolourable class of snarks, he ``kills three birds with one stone'', since it automatically implies the NZ $5$-flow, the ChNZ $5/2$-flow, and also the rich NZ $7$-flow. In a similar way, when someone shows a construction of some flow from the $1/2$-flow-pair, it is immediately proved for all classes implying $1/2$-flow-pair.

In Chapter \ref{ch:bounds}, we proved a lower bound on two-dimensional Chebyshev flow number of snarks. Although we are not aware of any infinite class meeting this bound, we described the infinite class of so-called \emph{Petersen book} graphs, which are not cubic, but their flow number corresponds to the lower bound for snarks of same order. Then we constructed snarks with order higher than the order of Petersen books, but with the flow number that is not greater than the flow number of Petersen books. This shows that the two-dimensional Chebyshev flow number can be arbitrarily near to $2$ for snarks.